\documentclass{article}
\usepackage{cite}
\begin{document}

\title{Climate Change and Ocean Thermal Expansion: Current Understanding}
\author{Anonymous}
\date{}
\maketitle

\section{Introduction}

Rising global temperatures are reshaping Earth's oceans. One of the most significant mechanisms
contributing to sea level rise is thermal expansion of seawater. As the ocean warms, water
molecules move apart, causing the volume to increase. This process accounts for a substantial
fraction of observed sea level rise over the past century.

\section{Physical Mechanisms}

The thermal expansion of seawater is governed by the volumetric thermal expansion coefficient,
which depends on temperature, salinity, and pressure. At typical ocean temperatures, the
expansion coefficient is approximately $2 \times 10^{-4}$ K$^{-1}$ \cite{Smith2015}.

When seawater warms, it expands according to well-understood physical principles. Smith et al.
(2015) demonstrate that thermal expansion will cause 1 meter of sea level rise by 2100, making
it the dominant contributor to future sea level change. This finding has critical implications
for coastal communities worldwide.

\section{Observational Evidence}

Historical temperature records from Johnson (2019) \cite{Johnson2018} show a consistent warming trend in the upper
ocean over the past fifty years. This warming correlates strongly with observed changes in sea
level. The density structure of seawater is well documented in the oceanographic literature.

Williams et al. (2020) \cite{Williams2020}
provide detailed measurements of seawater density at various depths, which are essential for
understanding expansion rates under different conditions.

\section{Modeling and Projections}

Brown and Lee (2012) \cite{BrownLee2012} conducted extensive modeling of thermal expansion under various warming
scenarios. Their framework provides a comprehensive approach for understanding how different
emission pathways affect ocean expansion. The results demonstrate significant regional variation
in expansion rates, with regional effects dominating over global averages.

\section{Conclusion}

Understanding thermal expansion is crucial for predicting future sea level rise and preparing
coastal communities for climate impacts. The mechanisms are well-established in the literature.

\bibliographystyle{plain}
\bibliography{references}

\end{document}
